\section{La contrazione muscolare}

\textbf{Teoria dello scorrimento dei filamenti}: durante la contrazione i filamenti sottili (actina) scorrono su quelli spessi (miosina) grazie al movimento dei \textbf{ponti trasversali}, senza che la lunghezza dei singoli filamenti cambi.
\begin{itemize}
  \item nel sarcomero rilasciato: banda A costante, zona H e banda I (emibanda I) più ampie;
  \item nel sarcomero contratto: la zona H e la banda I si riducono, mentre la banda A resta invariata, per effetto dello scorrimento dei filamenti di actina l'uno verso l'altro;
  \item il ponte trasversale muove il filamento di actina verso il centro del sarcomero.
\end{itemize}

% ---------------------------------------------------------------------------
\subsection{Proteine contrattili}

\rubrica{Filamento sottile}
Costituito da: F-actina (due catene), troponina (TnC, TnI, TnT) e tropomiosina.

\rubrica{Filamento spesso}
Costituito da molecole di miosina, ciascuna con due teste (\textbf{S1}, o \textbf{MHC}) e una coda (\textbf{S2}); sulla testa sono presenti i siti di legame per l'actina e per l'ATP (sede dell'attività ATPasica); le \textbf{catene leggere di miosina} (MLC) sono associate alle teste.
\begin{itemize}
  \item le teste delle molecole di miosina sono disposte in modo elicoidale, sfasate in senso lineare di 143\,\AA{} l'una rispetto alla successiva;
  \item la parte centrale del filamento spesso (\textbf{zona nuda}) è priva di ponti trasversali;
  \item la miosina può stabilire rapporti con l'actina per mezzo dei \textbf{ponti trasversali}.
\end{itemize}

% ---------------------------------------------------------------------------
\subsection{La placca neuromotoria}

Le fibre muscolari vengono portate in contrazione dalle \textbf{fibre nervose motorie}, tramite le giunzioni sinaptiche neuromuscolari o \textbf{placche motrici}. Il potenziale d'azione della fibra muscolare viene innescato dal potenziale postsinaptico delle placche motrici (\textbf{potenziale di placca}).
\begin{itemize}
  \item la durata totale del potenziale d'azione delle fibre muscolari è 2--5\,msec;
  \item il potenziale d'azione si propaga dal luogo d'insorgenza lungo l'intera fibra tramite le correnti elettrotoniche.
\end{itemize}

\rubrica{Componenti della placca}
Terminazione nervosa, sarcolemma, acetilcolina, recettore, mitocondri (schema: 1 = terminazione nervosa; 2 = sarcolemma; 3 = acetilcolina; 4 = recettore; 5 = mitocondri).

\rubrica{Struttura della giunzione neuromuscolare}
\begin{itemize}
  \item \textbf{cellula gliale} e \textbf{terminale presinaptico}, con vescicole sinaptiche contenenti \textbf{ACh};
  \item il potenziale d'azione in arrivo determina il rilascio di ACh nella \textbf{fessura sinaptica};
  \item sul sarcolemma della placca motrice sono presenti \textbf{pieghe giunzionali}, i siti recettoriali per l'acetilcolina e le molecole di \textbf{acetilcolinesterasi}.
\end{itemize}

% ---------------------------------------------------------------------------
\subsection{Tubuli trasversi, reticolo sarcoplasmatico e triade}

I \textbf{tubuli T} consentono il propagarsi della depolarizzazione dalla superficie verso l'interno della cellula muscolare. Il lume di ciascun tubulo contiene liquido extracellulare e comunica con l'esterno della fibra muscolare.
\begin{itemize}
  \item le due diramazioni di un tubulo T formano circonferenze che circondano ciascuna miofibrilla, interponendosi tra le due cisterne terminali;
  \item non esiste continuità morfologica tra tubuli T e cisterne terminali;
  \item l'insieme di due cisterne terminali e del tubulo T compreso tra di loro costituisce la \textbf{triade}.
\end{itemize}

% ---------------------------------------------------------------------------
\subsection{Accoppiamento eccitazione-contrazione}

\begin{enumerate}
  \item rilascio di ACh e legame ai recettori $\rightarrow$ arrivo di un potenziale d'azione nel tubulo T;
  \item il recettore diidropiridinico (\textbf{DHPR}), sensore del voltaggio sulla membrana del tubulo T, si attiva ed induce il rilascio di $Ca^{2+}$ dal reticolo sarcoplasmatico attraverso il recettore per la rianodina (\textbf{RyR\textsubscript{1}});
  \item esposizione dei siti attivi della G-actina (mascherati in condizioni di riposo dalla tropomiosina) e legame dei ponti crociati con l'actina;
  \item inizio della contrazione.
\end{enumerate}

\rubrica{Ruolo del calcio}
\begin{itemize}
  \item \textbf{stato di riposo}: i siti attivi della G-actina sono coperti dalla tropomiosina;
  \item \textbf{esposizione dei siti attivi}: il $Ca^{2+}$ si lega alla troponina, spostando la tropomiosina e scoprendo i siti attivi;
  \item \textbf{formazione dei ponti crociati}: la miosina si lega ai siti attivi dell'actina.
\end{itemize}

\rubrica{Regolazione del \texorpdfstring{$Ca^{2+}$}{Ca2+} nel sarcoplasma}
Al sarcolemma e ai tubuli T sono associati lo scambiatore $Na^+$/$Ca^{2+}$ e la pompa del $Ca^{2+}$ (\textbf{PMCA}); al reticolo sarcoplasmatico sono associati il recettore DHPR (accoppiato al RyR\textsubscript{1}) e la proteina di trasporto del calcio del reticolo sarcoplasmatico (\textbf{SERCA}), che pompa il $Ca^{2+}$ nuovamente nel reticolo terminando la contrazione.

% ---------------------------------------------------------------------------
\subsection{Il ciclo dei ponti}

\begin{enumerate}
  \item l'ATP si lega alla miosina, causandone il distacco dall'actina;
  \item la miosina idrolizza l'ATP, la testa della miosina ruota e si lega all'actina;
  \item segnale del $Ca^{2+}$, rilascio del $P_i$ e \textbf{colpo di frusta};
  \item la miosina rilascia l'ADP;
  \item stato di rigor.
\end{enumerate}

\rubrica{Legame dell'ATP e distacco dall'actina}
Il ciclo comincia con il legame di ATP alla testa della miosina:
$$M + ATP \rightarrow M \cdot ATP$$
La miosina contiene ATPasi, un enzima in grado di scindere il legame ad alta energia dell'ATP (idrolisi). Se la testa della miosina è legata all'actina ne consegue un rapido distacco dell'actina dal complesso actina-miosina-ATP:
$$AM + ATP \rightarrow M \cdot ATP$$

\rubrica{Idrolisi dell'ATP}
Un piegamento della regione del collo della miosina, con conseguente sollevamento della testa, accompagna l'idrolisi dell'ATP, con conseguente rilascio di energia. Parte di questa energia è assorbita dalla molecola di miosina, che raggiunge così lo stato ad alta energia ($M_0$), mentre i prodotti finali della reazione (ADP e $P_i$) rimangono fissati al sito ATPasico della miosina stessa ($M \cdot ADP \cdot P_i$):
$$M \cdot ATP \rightarrow M_0 \cdot ADP \cdot P_i$$

\rubrica{Legame con l'actina e rilascio del fosfato}
In seguito all'idrolisi dell'ATP, la miosina si lega all'actina ($A \sim M \cdot ADP \cdot P_i$); il legame è modulato dalla presenza di $Ca^{2+}$. Il legame forte della miosina all'actina è associato al rilascio del fosfato inorganico ($P_i$), che può così abbandonare il sito di legame dell'ATP della miosina:
$$A \cdot M \cdot ADP \cdot P_i \rightarrow A \cdot M \cdot ADP$$

\rubrica{Colpo di forza (\textit{power stroke})}
L'estensione della zona del collo della miosina, stirando di circa 10\,nm la parte elastica con produzione di forza, determina il cosiddetto \textbf{colpo di forza} o \textbf{power stroke}. L'accorciamento della lunghezza del collo della miosina causa lo scivolamento dei filamenti, spessi e sottili, gli uni sugli altri. L'ADP viene rilasciato dal complesso actina-miosina:
$$A \cdot M_0 \cdot ADP \rightarrow A \cdot M_0$$

\rubrica{Nuovo legame dell'ATP}
Per permettere alla miosina di attaccarsi ad un nuovo sito per l'actina e ripetere il ciclo, il legame tra actina e miosina deve essere spezzato: ciò avviene con il legame di una molecola di ATP alla testa della miosina. L'ATP si lega in modo rapido ed irreversibile al sito ATPasico della miosina, formando il complesso actina-miosina-ATP ($A \sim M \cdot ATP$). Il legame actina-miosina diventa debole, con conseguente rapido distacco dell'actina; in questa trasformazione la miosina passa da uno stato ad alta energia ($M_0$) ad uno stato a bassa energia ($M$):
$$A \cdot M_0 \rightarrow A \sim M \cdot ATP \rightarrow M \cdot ATP$$

\rubrica{Stato di rigor}
In assenza di ATP, il ciclo si arresta quando la testa della miosina è ritornata nello stato a bassa energia e actina e miosina restano fissate insieme, incapaci di staccarsi. Questa condizione è detta \textbf{stato di rigor} ($A \cdot M_0$).

\rubrica{Indipendenza dei ponti}
Il ciclo di ogni ponte si realizza indipendentemente dagli altri ponti; solo circa il 25--50\% dei ponti risultano attaccati e producono movimento in un dato istante.
